% mage-auto-heading: encoding-and-realization
% mage-auto-heading: cost-and-failure-separately
% mage-auto-heading: the-dynamic-model
% Authoritative LaTeX manuscript; Markdown is a generated compatibility projection.

% <!-- part-title: The Theory -->

% <!-- chapter-title: A Theory of MAGE -->

\chapter{A Theory of MAGE}

% <!-- point: the-pieces-this-chapter-assembles | Agentic engineering brings Modeling and Alignment together around two responsibilities: engineers encode consequential intent, and agents produce realizations that must fulfill it. -->

Agentic engineering must preserve engineering intent while delegating implementation. Modeling makes that intent available for reasoning; Alignment supplies independent evidence and enforcement. Within the software factory (\magevalue{chapter:software-factories}), engineers state consequential obligations and agents attempt to produce realizations that fulfill them.

% <!-- point: the-starting-point-is-the-master-equation | The encoding–realization decomposition separates adequate encoding from fulfillment; its chain rule does not itself establish the effect of an engineering intervention. -->

We can begin by separating those two responsibilities. Let \(E\) mean that the encoding captures the consequential intent and \(F\) that the resulting artifact fulfills the encoded obligations. For a representation \(R\), the probability of both is

\[
p_R = P(E \cap F \mid R) = \underbrace{P(E \mid R)}_{\substack{\text{Encoding: engineers state}\\\text{what must hold}}} \;\underbrace{P(F \mid E,R)}_{\substack{\text{Realization: agents fulfill}\\\text{those obligations}}}
\]

This \textbf{encoding–realization decomposition} is the chain rule; it requires no independence assumption. It tells us where to look for failure, but does not establish that a particular representation improves either factor. That requires evidence. Agents can help construct and revise the encoding, while engineers retain authority over the intent and agents retain freedom among acceptable realizations.

% <!-- point: the-probabilities-belong-to-a-production-system | The probabilities belong to a production system: the factory of Chapter 5, redrawn with each step of the realization path carrying the probability that it succeeds and the encoded property also supplying the tolerance the result is evaluated against. | terms: tolerance -->

\autoref{probability-in-the-software-factory} places the two terms on the software factory. The encoded property also supplies a tolerance against which the result can be checked independently of the agent that produced it.

\begin{figure}[htbp]
\centering
\includegraphics[width=\linewidth]{book/latex/assets/probability-in-the-software-factory-13c48ad51e.pdf}
\caption[Probability in the software factory --- encoding and realization, with the encoded property also supplying the tolerance]{\emph{Probability in the software factory.} Engineers encode consequential intent; agents fulfill the encoded obligations. The two factors condition on representation R, omitted here for legibility. Interpretation is a mechanism within realization. The encoded property also supplies an independently evaluable tolerance, so evidence and admission can govern it even when production fails.}
\label{probability-in-the-software-factory}
\end{figure}

\section{Encoding and Realization}\label{heading:encoding-and-realization}

% <!-- point: p-r-factors-into-encode-interpret-realize | Encoding and realization locate different interventions; faithful production through a representation is distinct from independently assessed task correctness. -->

The two terms distinguish failures in what the agent was asked to produce from failures in producing it. Holding the task population and production setup fixed, \(R\) names the representation choice; probabilities range over the encodings and realizations produced through it. The first question is whether the representation captures the consequential intent the engineer authorizes. The second is whether the resulting artifact fulfills those obligations.

Realization includes interpretation, planning, implementation, and correction. An agent may initially misunderstand an obligation and then correct its work after a check or review exposes the error. Conversely, it may infer an obligation that the encoding omitted. The event \(E\cap F\) therefore describes \emph{faithful production} through the representation; observed task correctness must be assessed separately.

% <!-- point: the-encoding-term-is-also-a-human-channel | A representation can improve P(E | R) by helping the engineer express intent correctly even if conditional realization is unchanged; Modeling improves the human channel as well as the agent channel. -->

The first term is easy to overlook in an agentic setting. Representation is not valuable only because a machine may interpret a structured model more reliably than prose. It can also help the engineer express the intended property correctly in the first place. Different representations make different distinctions visible, force different decisions to be made explicitly, and make omissions easier or harder to notice. A representation can therefore improve \(P(E \mid R)\) even if it leaves \(P(F \mid E,R)\) unchanged. This is one reason Modeling is an engineering activity rather than merely a technique for supplying context to an agent.

% <!-- point: capable-agents-can-raise-the-encoding-term-through-feedback | Feedback from a failed realization can help diagnose an encoding problem; capable agents may improve encoding indirectly when that diagnosis supports revising intent. -->

Representation is not the only route to that term. A realization that misses the engineer\textquotesingle s meaning can provide evidence about the encoding, although the failure may instead lie in realization. The comprehension amplification of \magevalue{sec:the-dynamics-of-mage} can help the engineer diagnose the cause and revise the encoding when needed. In that way, capable agents can contribute indirectly to \(P(E \mid R)\).

% <!-- point: the-two-factory-techniques-act-on-different-factors | Process design changes the conditions and cost of realization. Explicit tolerances help encode intent and make selected properties available for independent evaluation. | terms: tolerance -->

Process design changes the conditions under which agents realize intent: context, tools, decomposition, routing, permissions, and other features of the harness can change \(P(F\mid E,R)\), its cost, or both. A better agent can improve the same term. Explicit tolerances help engineers state the intended property and make it available for independent evaluation. A check can then compare the realization with the represented tolerance, and a gate can reject work that violates it. This is the distinction between helping production succeed and governing what happens when it fails.

% <!-- point: stronger-reasoners-can-recover-meaning | Better reasoners may reconstruct meaning from a weak representation, but some lost distinctions require a different representation or additional evidence. -->

A stronger reasoner can sometimes recover meaning that a weaker one misses. A terse requirement may be enough for an agent that knows the domain and can inspect the surrounding system. Better retrieval, tools, and feedback can also help it discover what the representation leaves implicit. These possibilities affect the cost and reliability of realization; they are reasons to compare richer representations with more capable inference rather than assume that either must win.

There are limits to what can be recovered from the available evidence. The entitlement failure from \magevalue{chapter:modeling} illustrates the problem. A scalar representation read each group's unlimited flag and still returned a single finite sum: its output type had no place for the distinction. A reasoner using only that output could not distinguish unlimited entitlement from the finite sum it received. Repair required revisiting the representation or consulting information outside it.

% <!-- point: reliable-realization-does-not-repair-wrong-encoding | With encoding quality fixed, perfectly reliable fulfillment leaves faithful production bounded by the encoding term; this does not imply complete specification. -->

Even perfectly reliable realization cannot, by itself, establish that the encoded obligations were right. Holding encoding quality fixed, as conditional realization approaches certainty, faithful production approaches \(P(E\mid R)\). The practical consequence is to examine the encoding when better agents keep faithfully producing the wrong thing.

% <!-- point: the-bound-does-not-argue-for-complete-specification | Useful freedom lets agents choose satisfactory realizations while engineers specify consequential obligations; complete specification is not the objective. -->

The appropriate degree of specification depends on the work. Some transformations leave little useful freedom, and the cost of ambiguity may outweigh any benefit from delegation. Elsewhere, specifying nearly every realization decision in advance recreates much of the specification cost discussed in \magevalue{chapter:software-factories}. It also sacrifices a major advantage of commodity intelligence: agents can explore implementation choices that engineers need not make themselves. The engineer can specify consequential obligations, leave other choices open, and let the agent find a satisfactory realization. The economic value of Modeling and Alignment therefore depends partly on how much useful freedom the environment can preserve while maintaining adequate control.

\section{Lifecycle Cost and Engineering Value}\label{heading:cost-and-failure-separately}

% <!-- point: lifecycle-cost-counts-correct-delivery | Cost per correct admitted change counts construction, all attempts and upkeep over comparable requested work and a declared horizon. -->

The engineering question is what useful delivery costs. A cheap attempt may leave an engineer with hours of repair; an expensive model may support hundreds of later changes. To compare those arrangements, we need to count the work that went into delivery, including the work that never produced an acceptable result.

For \(N\) requested changes, let \(D_N\) count the changes independently verified as correct and admitted within a declared evaluation horizon. Each requested change can include several attempts. The total cost and the cost per correct admitted change are

\[
C_N=C_{\text{build}}+\sum_{i=1}^{N}\bigl(C_{\text{run},i}+C_{\text{carry},i}\bigr),
\qquad c_N=\frac{C_N}{D_N},\quad D_N>0.
\]

Construction covers the initial investment in the engineering environment. Running costs include production, checking, human review, and repair across all attempts; carrying costs cover the upkeep needed to keep the environment useful. A failed or abandoned task still contributes its costs to the numerator. If nothing correct is delivered, report that result and the expenditure rather than assigning a cost per delivery. Shared costs need an explicit allocation, and human effort and machine expenditure need declared valuations before they can be added. Elapsed time remains a separate outcome.

% <!-- point: cost-comparisons-retain-work-quality-and-horizon | Comparable requested work, quality and horizon prevent an attractive cost ratio from concealing rejected or unfinished work. -->

This accounting does not assume that tasks cost the same, that retries are independent, or that the environment stands still. It does require comparable work. An engineering process that accepts only easy changes may have an attractive \(c_N\) while leaving important requests unfinished. Comparisons should therefore hold the requested work and evaluation horizon fixed, report what was delivered and left unresolved, and measure escaped failures and false rejection alongside cost. Repeated attempts often change the specification, tools, or next plan; observed delivery costs can count that learning without pretending that every retry starts over under identical conditions.

% <!-- point: assurance-without-modeling-can-be-expensive | Validation, false rejection, and repair consume running cost even when the admission mechanism’s protection stays unchanged. -->

Admission adds the cost of gathering evidence and acting on it. A gate can provide the same protection while one representation causes agents to collide with it much more often than another. Each rejected candidate may require further inference, diagnosis, or human attention. These costs belong in \(C_{\text{run},i}\), whether the rejection catches a real defect or mistakenly turns away correct work.

% <!-- point: carrying-cost-completes-the-comparison | A representation earns its construction and upkeep through useful work enabled or later costs avoided, compared with inference over less explicit knowledge. -->

Carrying cost completes the comparison. The factories in \magevalue{chapter:software-factories} preserve engineering knowledge in models, instructions, policies, skills, tests, and context systems. All require some effort to remain useful as the software changes. A structured model earns its construction and upkeep when it enables valuable work or reduces later reconstruction, realization, failure, or assurance costs enough over its useful life. For a simple or infrequent task, a natural-language description may already do the job at very little cost.

The alternative is often to buy inference when the question arises. A capable agent may reconstruct a dependency from source cheaply enough that a maintained graph would not repay its cost. For another task family, reconstructing that same dependency hundreds of times may be much more expensive than keeping it explicit. The comparison depends on the work, the available reasoner, the assurance required, and the expected reuse of the representation. A model can also repay its cost in a single use when the question is sufficiently consequential.

% <!-- point: vendor-factory-tests-representation-versus-inference | A packaged agent system may substitute capable inference and generic machinery for some explicit structure; compare what it delivers, costs, and leaves for later work. -->

This gives packaged agent systems a place in the theory. A capable reasoner and generic supervisory machinery may economically replace some structure that a weaker reasoner would need. How much can they replace, and for which obligations? We can compare what those arrangements deliver, what escapes their assurance, and what they cost, while also asking what engineering knowledge survives for the next task.

% <!-- point: human-attention-belongs-in-lifecycle-cost | Human review and intervention consume engineering resources even when machine execution is inexpensive. -->

Human attention is part of this accounting. A production system that generates changes cheaply but requires an expert to reconstruct each one has moved work onto that expert. The relevant cost includes the time needed to interpret the evidence, diagnose failures, and decide whether the result is acceptable. The following inset explains why increasing output can make this cost decisive.

% <!-- box-family: inset -->

% <!-- inset-domain: HUMAN FACTORS -->

% mage-box-prefix: Inset —

\begin{bookbox}{inset}{Human Attention Is Not a Free Control}
\label{inset-human-attention-not-a-free-control}
Moving work from human execution to machine execution does not eliminate the cost of human attention if people must inspect the resulting stream. Software engineering has long relied on one engineer reviewing another\textquotesingle s work, but code review is imperfect. Studies document the improvements reviewers prompt and associate review coverage and participation with post-release quality.\citep{beller2014modern, mcintosh2014impact} At Microsoft, reviews spanning more files had a lower proportion of useful feedback, and developers reported that such reviews required more time and effort.\citep{bosu2015useful} As agents assume more implementation work, software engineering faces a problem familiar to industries that automated earlier. Bainbridge\textquotesingle s \emph{Ironies of Automation} described how automation can leave human operators responsible for monitoring rare failures and intervening when something goes wrong.\citep{bainbridge1983}

As agentic output grows, proportional human inspection becomes a bottleneck. Where recurring obligations can be represented explicitly and checked or enforced adequately, reusable mechanisms can reduce the inspection required for each change. Engineers can then concentrate their attention on ambiguous evidence, novel failures, tradeoffs, exceptions, and decisions whose consequences require expert judgment.
\end{bookbox}

% <!-- point: human-attention-is-a-conserved-resource | Human attention is an engineering resource. -->

Human attention is an engineering resource.

% <!-- point: factory-value-requires-more-than-correct-delivery | Correct delivery is a constraint, not a universal valuation of outcomes; upkeep, consequences and intended use determine whether the arrangement pays. -->

A low cost per correct change still does not establish that the changes were worth making. Product work may be valued for capabilities delivered to customers; security work, for risk retired; accessibility work, for documents made accessible and remediation effort avoided. The accounting names the expenditure, while the engineering decision must also consider the value of the result and the consequences of failure. A repair cost already charged to running the engineering process should not be counted again as a separate loss.

Nor does an idle engineering environment necessarily become free. It may retain operating costs, and its engineering assets may need maintenance as dependencies, requirements, and tools change. Some structures remain useful for years with little attention; others lose fit quickly. This is why the return on an environment depends on its useful life and carrying cost, rather than on a one-time demonstration that it can produce correct software.

\section{The Inherited Engineering Environment}\label{heading:the-dynamic-model}

% <!-- point: se-needs-a-model-where-each-realization-changes-the-next | Software engineering extends the episode model: each realization becomes structure inherited by the next. -->

Software engineering extends the model beyond one episode: each realization changes the conditions inherited by the next. Winters et al. characterize software engineering as ``programming integrated over time'': software is developed, modified, and maintained as both the system and its environment change.\citep{winters2020} One realization therefore becomes part of the architecture, conventions, representations, tests, and other structure inherited by subsequent work. A good change can make later work easier to reason about; a locally acceptable but structurally poor change can make it harder.

% <!-- point: independence-omits-the-inherited-environment | Each engineering episode inherits an environment and may change the conditions of later work; the dynamic relation is schematic, not a fitted law. -->

The lifecycle account totals the work and its costs. It does not yet explain how earlier work changes the conditions for later production. Let \(T_i\) be the task at step \(i\), \(M\) the foundation model, \(R_i\) the representation supplied for that task, and \(H\) the harness and process. Write \(\mathcal{G}_i\) for the governed engineering environment inherited at that step and \(a_i\) for the work and resulting changes that the episode leaves behind:

\[
p_i=p(T_i,M,R_i,H,\mathcal{G}_i),
\qquad \mathcal{G}_{i+1}=f(\mathcal{G}_i,a_i).
\]

Here \(p_i\) is the probability of faithful production for task \(T_i\), extending the earlier encoding–realization account to include the inherited environment. The second says that the episode can change that environment. A failed attempt may leave a useful diagnosis or a new control even when it produces no admitted change. The notation holds the model and harness fixed to expose inheritance; they too can change across a longer sequence. These are schematic relationships, not calibrated laws or a claim that every episode improves the environment.

% <!-- point: coupling-creates-path-dependence-decay-self-reinforcing | The coupling creates path dependence: poor structure makes good realization harder, and decay can become self-reinforcing. -->

This coupling creates path dependence. A locally acceptable change can increase architectural irregularity, obscure structure, or introduce inconsistent conventions, leaving later agents a worse reasoning surface. A change that improves representations, architecture, constraints, or other engineering capital can instead make later work easier to reason about and govern.

% <!-- point: theory-of-mage-informal-statement | Agentic capacity acts through an inherited environment, shaping durable throughput, defect escape and human-attention burden. -->

Engineering outcomes depend on the interaction between agentic capacity and the environment in which that capacity operates. \textbf{Agentic capacity} supplies model capability, concurrency, and raw implementation rate. The \textbf{governed engineering environment} supplies the representations, constraints, evidence, and enforcement through which that capacity acts. Together they shape durable throughput, defect escape, and human attention consumed per unit of useful output. \autoref{mage-dynamic-model} draws the interaction and its feedback dynamics.

\begin{figure}[htbp]
\centering
\includegraphics[width=\linewidth]{book/latex/assets/mage-dynamic-model-62588a5597.pdf}
\caption{\emph{The MAGE Dynamic Model.} Agentic capacity acts through a governed engineering environment. Mismatch creates governance pressure; diagnosis and adaptation can improve the environment inherited by later work. The balancing loop seeks fit, while capability amplification reopens the frontier as stronger environments support more ambitious work. Prior engineering knowledge seeds the system before failure occurs.}
\label{mage-dynamic-model}
\end{figure}

% <!-- point: diagnosis-and-authority-connect-pressure-to-adaptation | Pressure prompts useful adaptation only when someone can diagnose the structural cause and change the shared environment. -->

The feedback requires engineering judgment. Repeated failures can expose an inadequate model or missing control, but they can also provoke another local patch. Someone must distinguish the two and have a path to change the shared environment. Adaptation may add a representation, strengthen an admission mechanism, reconcile conflicting rules, or retire a control that no longer earns its cost. Prior engineering knowledge can also seed that environment before a failure occurs; MAGE does not require learning every lesson through an incident.

% <!-- point: four-properties-of-an-effective-environment | Useful representations, adequate coverage, coherence and economy describe environment fit independently of machinery counts. -->

% <!-- print-keep-with-next: 8 -->

More machinery does not mean a better environment. Ten thousand lines of stale checks are worse than one hundred lines that govern the obligations that matter. Four properties help us ask whether the environment fits the work:

\begin{enumerate}
\tightlist
\item \textbf{Useful representations.} The information needed for the work is explicit at an appropriate level of abstraction.
\item \textbf{Adequate coverage.} Important obligations have appropriate evidence and enforcement.
\item \textbf{Coherence.} Models and mechanisms reinforce one another rather than contradicting or needlessly duplicating.
\item \textbf{Economy.} Their future value justifies their maintenance, friction, and context cost.
\end{enumerate}

Coverage does not require mechanizing every obligation. Some belong in a compiler rule or gate; others require statistical evidence, semantic evaluation, or human judgment. A strong environment makes clear where judgment belongs and where repeatedly buying the same judgment is wasteful.

% <!-- point: not-a-maturity-model-improve-the-limiting-part | The properties are dimensions for diagnosing what limits the work, not maturity stages; the appropriate response may improve representations or assurance, or reconcile and retire controls. -->

These properties describe different ways an environment can support the work, rather than stages of a maturity model.\citep{bollinger1991critical} The appropriate investment depends on what is limiting progress. Repeated reconstruction of important state suggests a representation problem; escaped failures call for examining evidence and enforcement. When existing controls impose more friction than protection, reconciliation or retirement may improve the environment more than adding another mechanism.

% <!-- point: definitions-of-the-three-outcomes | Durable throughput, defect escape and human-attention burden are outcomes; their candidate observables require explicit measurement choices. -->

A theory of agentic engineering should distinguish activity from progress. \textbf{Durable throughput} is useful system progress that survives rework, integration, and operation. \textbf{Defect escape} records consequential failures that reach later lifecycle stages or production. \textbf{Human-attention burden} measures the scarce engineering judgment consumed per unit of durable progress. \autoref{mage-outcomes} gives candidate observables; \magevalue{sec:research-agenda} returns to how comparisons can distinguish the effects of representation, enforcement, and inherited capital.

{\footnotesize
\begin{longtable}[]{@{}
  >{\raggedright\arraybackslash}p{(\linewidth - 4\tabcolsep) * \real{0.3333}}
  >{\raggedright\arraybackslash}p{(\linewidth - 4\tabcolsep) * \real{0.3333}}
  >{\raggedright\arraybackslash}p{(\linewidth - 4\tabcolsep) * \real{0.3333}}@{}}
\caption[The three MAGE outcomes and their observables]{\emph{Outcomes of the theory.} Raw implementation activity is an input. The theory is interested in how much survives as durable progress, what failures escape, and how much scarce human judgment that progress consumes.}\label{mage-outcomes}\tabularnewline
\toprule\noalign{}
\begin{minipage}[b]{\linewidth}\raggedright
Outcome
\end{minipage} & \begin{minipage}[b]{\linewidth}\raggedright
What it measures
\end{minipage} & \begin{minipage}[b]{\linewidth}\raggedright
Candidate observables
\end{minipage} \\
\midrule\noalign{}
\endfirsthead
\toprule\noalign{}
\begin{minipage}[b]{\linewidth}\raggedright
Outcome
\end{minipage} & \begin{minipage}[b]{\linewidth}\raggedright
What it measures
\end{minipage} & \begin{minipage}[b]{\linewidth}\raggedright
Candidate observables
\end{minipage} \\
\midrule\noalign{}
\endhead
\bottomrule\noalign{}
\endlastfoot
\textbf{Durable throughput} & Raw implementation converted into lasting system progress & landed changes net of rework; churn share; the slope of sustained throughput \\
\textbf{Defect escape} & Relevant failures reaching later lifecycle stages or production & escaped-defect rate; reopen rate; silent-corruption escape \\
\textbf{Human-attention burden} & Human judgment required per unit of durable output & interventions per landed change; review time; conversion effort \\
\end{longtable}
}

% <!-- point: activity-counts-are-not-outcome | Easy-to-count quantities (commits, lines, control files, support ratio) measure activity or investment, not outcome. -->

Several easy-to-count quantities measure activity or investment rather than outcome: commits, lines changed, control files, model-sync events, and support ratio. They can help explain what an environment is doing; they do not establish that it is producing a return.\magefootnote{activity-data}{The activity-versus-outcome gap also appears in industry data. LinearB\textquotesingle s 2026 benchmark covers 83,000 developers across 253 organizations and compares merge rates from June 2025 to May 2026.\citep{linearb2026} Among the heaviest AI users, merge rates roughly doubled while pull-request yield---the share of opened requests merged within 30 days---fell. The association does not establish that AI caused either change, and neither measure establishes durable progress. An industry analyst reads the same period as rising AI spend without proportionate shipping velocity, the difficulty having moved from comprehension to confidence.\citep{riggins2026}}

% <!-- point: three-anchors-open-four-propositions | The production identity, lifecycle accounting and inherited environment motivate four conditional propositions requiring stronger empirical tests. -->

The production equation separates two responsibilities, lifecycle accounting counts the resources spent, and the dynamic relation makes inheritance explicit. None supplies a measured effect size. Together they support a more focused question: under what conditions do representations, independent enforcement, and inherited engineering judgment turn additional agentic capacity into useful progress? The next section develops four propositions that evidence could challenge.
